\documentclass[10pt,a4paper]{amsart}
\usepackage[margin=19mm]{geometry}
\usepackage{amsmath,amssymb,mathtools}
\usepackage[scr=rsfs,cal=euler]{mathalfa}
\usepackage{xfrac}
\usepackage[unicode,hidelinks,bookmarks=false]{hyperref}
\newcommand{\ie}{i.e.\ }
\newcommand{\cf}{cf.\ }
\newcommand{\resp}{respectively\ }
\newcommand{\Subsection}{\subsection}
\providecommand{\nobreakdash}{\nobreak\hbox{-}}
\setlength{\emergencystretch}{2em}
\multlinegap0pt
\usepackage{amssymb}
\usepackage{dsfont}
\newtheorem{lemma}{Lemma}
\def\d {{\rm d}}
\title{Traveling waves for SIR model on two-dimensional lattice}
\author{Ran Zhang, Shunchang Su, Xue Ren}
\date{Source: arXiv:2602.16172v1 (CC BY 4.0)}
\begin{document}
\maketitle

\noindent\emph{Source and licence.} Traveling waves for SIR model on two-dimensional lattice. Version arXiv:2602.16172v1 (CC BY 4.0), distributed by arXiv under the Creative Commons Attribution 4.0 International licence, http://creativecommons.org/licenses/by/4.0/, as recorded in the arXiv licence metadata for this exact version and shown on the abstract page https://arxiv.org/abs/2602.16172v1. Full licence notice: \texttt{licenses/arxiv-2602.16172-CC-BY-4.0.txt}.\par\smallskip

\noindent\textit{Editorial scope.} Counted unit: the article's lemma 10 (source0.tex lines 803--805). The article's complete original proof of it is delivered with it (source0.tex lines 806--869, one \texttt{proof} environment of 64 lines). No mathematics is written, repaired or re-derived: the statement and the proof are the article's own lines, in the article's own notation, and no retained line is substituted or re-flowed.  Retained uncounted as the source-local context the proof needs: lines 206--215.  Nothing was omitted from any retained span, so the package carries no omission record.  No retained line contains a citation, so the wrapper carries no bibliography and no bibliography entry.  No retained line contains a figure environment, an image-inclusion command or drawing code, so no image is delivered and the package declares no asset dependency.  Editorial formatting only, in addition: an AMS \texttt{amsart} preamble that restates the article's own declarations and macros used by the retained lines, namely the environments \texttt{lemma} on separate counters, together with the standard packages those lines need.  The retained environments print in this document as Lemma 1 in order of appearance, whereas the article prints the same items with the numbering of its own full document.  The complete native source of the article as distributed in the arXiv e-print for this exact version is bundled unchanged as \texttt{sources/194854/source0.tex}.
\begin{equation}
\label{WaveEqu}\left\{
\begin{array}{l}
\vspace{2mm}
\displaystyle   c S'(\xi) =
d_1\mathfrak{J}[S](\xi) + \Lambda - \frac{\beta S(\xi)I(\xi)}{1+\alpha I(\xi)} - \mu_1 S(\xi),\\
\displaystyle   c I'(\xi) =
d_2\mathfrak{J}[I](\xi) + \frac{\beta S(\xi)I(\xi)}{1+\alpha I(\xi)} - \mu_2 I(\xi),
\end{array}\right.
\end{equation}

\begin{lemma}\label{10}
The functions \(\displaystyle\frac{I(\xi\pm\sin\theta)}{I(\xi)}\), \(\displaystyle\frac{I(\xi\pm\cos\theta)}{I(\xi)}\) and \(\displaystyle\frac{I'(\xi)}{I(\xi)}\) are bounded in \(\mathbb{R}\).
\end{lemma}

\begin{proof}
Due to the uncertainty of the positive and negative values of  \(\sin\theta\) and \(\cos\theta\), we will categorize them into four situations. We assume  \(\theta\neq \displaystyle\frac{\pi}{2}+2k\pi~~(k\in  Z)\), firstly \(\sin\theta>0~,~\cos\theta>0\), secondly \(\sin\theta>0~,~\cos\theta<0\), thirdly \(\sin\theta<0~,~\cos\theta<0\), and finally \(\sin\theta<0~,~\cos\theta>0\). The proof ideas for the above four situations are the same. We will use the first situation as an example to demonstrate, and leave the remaining three situations for readers to prove themselves.

 From the second equation in the traveling wave system (\ref{WaveEqu}), we get
\[
cI'(\xi)+(4d_2+\mu_2)I(\xi)=d_2I(\xi + \sin\theta)+d_2I(\xi -\sin\theta)+d_2I(\xi +\cos\theta)+d_2I(\xi -\cos\theta)+
\frac{\beta S(\xi)I(\xi)}{1+\alpha I(\xi)}>0.
\]
Denote \(W(\xi)=e^{\nu\xi}I(\xi)\), where \(\nu=(4d_2+\mu_2)/c\). As a result
\[
cW'(\xi)=e^{\nu\xi}(cI'(\xi)+(4d_2+\mu_2)I(\xi))>0,
\]
thus \(W(\xi)\) is increasing on \(\xi\). Then, \(W(\xi - \sin\theta)<W(\xi)\), this means,
\[
\frac{I(\xi - \sin\theta)}{I(\xi)}<e^{\nu} ~~~ \forall ~ \xi\in\mathbb{R}.
\]
Note that
\begin{equation}\label{3.1}
\begin{aligned}
\left[e^{\nu\xi}I(\xi)\right]' &= \frac{1}{c}e^{\nu\xi}\left[d_2I(\xi + \sin\theta)+d_2I(\xi -\sin\theta)+d_2I(\xi +\cos\theta)+d_2I(\xi -\cos\theta)+
\frac{\beta S(\xi)I(\xi)}{1+\alpha I(\xi)}\right]\\
&>\frac{d_2}{c}e^{\nu\xi}I(\xi + \sin\theta).
\end{aligned}
\end{equation}
%Integrating \ref{3.1} over \([\xi, \xi + \sin\theta]\) and using the fact that \(e^{\nu\xi}I(\xi)\) is increasing, we have
Integrating \ref{3.1} on \([\xi, \xi + \sin\theta]\) and utilizing the fact that \(e^{\nu\xi}I(\xi)\) increases, we obtain
\begin{align*}
e^{\nu(\xi + \sin\theta)}I(\xi +  \sin\theta)&>e^{\nu\xi}I(\xi)+\frac{d_2}{c}\int_{\xi}^{\xi +  \sin\theta}e^{\nu s}I(s +  \sin\theta)ds\\
&>e^{\nu\xi}I(\xi)+\frac{d_2}{c}\int_{\xi}^{\xi +  \sin\theta}e^{\nu(\xi +  \sin\theta)}I(\xi +  \sin\theta)e^{-\nu\sin\theta}ds\\
&=e^{\nu\xi}\left[I(\xi)+\frac{d_2}{c}I(\xi +  \sin\theta) \sin\theta\right].
\end{align*}
By (\ref{3.1}), we have
\begin{equation}\label{3.2}
\begin{aligned}
\left[e^{\nu\xi}I(\xi)\right]'&>\left(\frac{d_2}{c}\right)^2e^{-2\nu\sin\theta}e^{\nu(\xi + \sin\theta)}I(\xi + \sin\theta)\sin\theta.
\end{aligned}
\end{equation}
Integrating (\ref{3.2}) over \(\left[\xi - \displaystyle \frac{\sin\theta}{2}, \xi\right]\) yields
\begin{align*}
e^{\nu\xi}I(\xi)&>\left(\displaystyle \frac{d_2}{c}\right)^2e^{-2\nu\sin\theta}\sin\theta\int_{\xi - \frac{\sin\theta}{2}}^{\xi}e^{\nu(s + \sin\theta)}I(s + \sin\theta)ds\\
 &>\left(\frac{d_2}{c}\right)^2\frac{e^{-2\nu\sin\theta}\sin^{2}\theta }{2}e^{\nu(\xi + \frac{\sin\theta}{2})}I\left(\xi + \frac{\sin\theta}{2}\right),
\end{align*}
that is
\[
\frac{I\left(\xi + \frac{\sin\theta}{2}\right)}{I(\xi)}<2\left(\frac{c}{d_2}\right)^2e^{\frac{3}{2}\nu\sin\theta}\sin^{2}\theta ~~\text{ for all } \xi \in \mathbb{R}.
\]
Similarly, integrating (\ref{3.2}) over \(\left[\xi, \xi + \frac{\sin\theta}{2}\right]\), we have
\[
\frac{I(\xi + 1)}{I\left(\xi + \frac{\sin\theta}{2}\right)}<2\left(\frac{c}{d_2}\right)^2e^{\frac{3}{2}\nu\sin\theta}\sin^{2}\theta~~ \text{ for all } \xi \in \mathbb{R}.
\]
Thus,
\[
\frac{I(\xi + \sin\theta)}{I(\xi)}=\frac{I\left(\xi + \frac{\sin\theta}{2}\right)}{I(\xi)}\frac{I(\xi + \sin\theta)}{I\left(\xi + \frac{\sin\theta}{2}\right)}<4\left(\frac{c}{d_2}\right)^4e^{3\nu\sin\theta}\sin^{4}\theta~~ \text{ for all } \xi \in \mathbb{R}.
\]
Similarly, the proof for $\displaystyle\frac{I(\xi \pm \cos\theta)}{I(\xi)}$ is similar.

By the second equation of (\ref{WaveEqu}), it follows that
\begin{align*}
cI'(\xi)&=d_{2}I(\xi + \sin\theta)+d_{2}I(\xi - \sin\theta)+d_{2}I(\xi + \cos\theta)+d_{2}I(\xi - \cos\theta)+\frac{\beta S(\xi){I(\xi)}}{1+\alpha I(\xi)}-(4d_2+\mu_2)I(\xi)\\
&\leq\d_{2}\frac{I(\xi + \sin\theta)}{I(\xi)}+d_{2}\frac{I(\xi - \sin\theta)}{I(\xi)}+d_{2}\frac{I(\xi + \cos\theta)}{I(\xi)}+d_{2}\frac{I(\xi - \cos\theta)}{I(\xi)}+\beta S(\xi)-(4d_2+\mu_2)\\
&\leq\d_{2}\frac{I(\xi + \sin\theta)}{I(\xi)}+d_{2}\frac{I(\xi - \sin\theta)}{I(\xi)}+d_{2}\frac{I(\xi + \cos\theta)}{I(\xi)}+d_{2}\frac{I(\xi - \cos\theta)}{I(\xi)}+\beta S_{0}-(4d_2+\mu_2),
\end{align*}
which gives us \(\displaystyle\frac{I'(\xi)}{I(\xi)}\) is bounded in \(\mathbb{R}\). This proof is over.
\end{proof}

\end{document}
